% ------------------------------------------------------------------------------
% Este fichero es parte de la plantilla LaTeX para la realización de Proyectos
% Final de Grado, protegido bajo los términos de la licencia GFDL.
% Para más información, la licencia completa viene incluida en el
% fichero fdl-1.3.tex

% Copyright (C) 2012 SPI-FM. Universidad de Cádiz
% ------------------------------------------------------------------------------

Este capítulo trata sobre todos los aspectos relacionados con la implementación del sistema en código, haciendo uso de un determinado entorno tecnológico.



\section{Arquitectura BI}
Para la contrucción de la arquitectura BI, hemos seguido un proceso ETL (\textit{Extract, Transform, Load}) que abarca desde la extracción de los datos en crudo mediante Api-Football hasta la carga en la base de datos PostgreSQL. Asi que, estructuraremos la explicación en tres fases: Extracción, Transformación y Carga, explicando la organización de directorios, archivos creados y conexiones entre las diferentes fases.
\subsection{Extracción}

En esta primera etapa se ha realizado la recopilación de datos de los últimos 10 años de la Liga Española de Fútbol procedente de una API externa especializada en fútbol, llamada API-Football, la cual contaba con una buena cantidad de datos sobre partidos, temporadas y rendimiento de equipos y jugadores. Para esta recopilación se han desarrollado varios \textit{scripts} en Python ubicados en el directorio DATOS encargados de hacer las diferentes peticiones a la API correspondientes a los datos que necesitaban, y a su vez almacenar los datos en formato JSON separados por temporada.

No se ha realizado transformación directa en estos \textit{scripts}, ya que se ha decidido que guardar los datos brutos era una buena práctica en busca de consultar los datos en busca de posibles inconsistencias, fallos o falta de algún dato. Ademas de poder repetir el proceso ETL sin necesidad de consultar cada vez a la API.

Los datos que se han extraído en esta etapa han sido los siguientes:

\renewcommand{\arraystretch}{1.2}

\begin{longtable}{|p{3.5cm}|p{4.9cm}|p{7cm}|}
\caption{Tipos de datos extraídos durante el proceso ETL}
\label{tab:tipos_datos_etl} \\

\hline
\textbf{Tipo de dato} & \textbf{Directorio de salida} & \textbf{Descripción} \\ 
\hline
\endfirsthead

\hline
\textbf{Tipo de dato} & \textbf{Directorio de salida} & \textbf{Descripción} \\ 
\hline
\endhead

\hline
\multicolumn{3}{r}{\textit{Continúa en la siguiente página}} \\
\hline
\endfoot

\hline
\endlastfoot

Equipos &
\texttt{DATOS\slash equipos\_base\slash} &
Información básica de los equipos participantes. \\ \hline

Partidos &
\texttt{DATOS\slash partidos\_base\slash} &
Calendario, jornadas, fechas, equipos locales, visitantes y resultados. \\ \hline

Estadísticas de partidos &
\texttt{DATOS\slash partidos\_stats\slash} &
Estadísticas agregadas de los equipos en cada partido. \\ \hline

Eventos de partidos &
\texttt{DATOS\slash partidos\_eventos\slash} &
Goles, tarjetas, sustituciones y otros eventos ocurridos durante los encuentros. \\ \hline

Jugadores por partido &
\texttt{DATOS\slash jugadores\_partido\slash} &
Rendimiento individual de cada jugador en cada partido. \\ \hline

Jugadores por temporada &
\texttt{DATOS\slash jugadores\_temporada\slash} &
Perfiles y estadísticas acumuladas de los jugadores por temporada. \\ \hline

Estadísticas de equipos por temporada &
\texttt{DATOS\slash equipo\_stats\_temporada\slash} &
Datos agregados del rendimiento de cada equipo durante una temporada. \\ \hline

\end{longtable}

Para esta extracción se han realizado los siguientes \textit{scripts}, los cuales explicaremos su funcionamiento:

\begin{longtable}{|p{5cm}|p{9cm}|}
\caption{Scripts desarrollados para la extracción de datos}
\label{tab:scripts_extraccion} \\

\hline
\textbf{Script} & \textbf{Función} \\ \hline
\endfirsthead

\hline
\textbf{Script} & \textbf{Función} \\ \hline
\endhead

\hline
\multicolumn{2}{r}{\textit{Continúa en la siguiente página}} \\ \hline
\endfoot

\hline
\endlastfoot

\path{extraer_equipos.py} &
Consulta la API para obtener la información de los equipos participantes en LaLiga para cada temporada. \\ \hline

\path{extraer_partidos.py} &
Extrae el calendario completo de partidos correspondiente a cada temporada. \\ \hline

\path{extraer_stats_partidos.py} &
Obtiene las estadísticas agregadas de cada equipo en los distintos partidos disputados. \\ \hline

\path{extraer_eventos_partidos.py} &
Recupera los eventos ocurridos durante los partidos, como goles, tarjetas, sustituciones y otras incidencias. \\ \hline

\path{extraer_jugadores_partido.py} &
Extrae las estadísticas individuales de cada jugador para todos los partidos disputados. \\ \hline

\path{extraer_jugadores_temporadas.py} &
Obtiene las estadísticas acumuladas y el rendimiento de los jugadores en cada temporada. \\ \hline

\path{extraer_equipos_temporada.py} &
Extrae las estadísticas agregadas de cada equipo al finalizar cada jornada de la temporada. \\ \hline

\path{extraer_jugadores_faltantes.py} &
Recupera la información de aquellos jugadores que no fueron obtenidos durante la extracción inicial. \\ \hline

\end{longtable}



Aunque existen distintos ficheros para extraer equipos, partidos, clasificaciones, estad\'isticas, eventos o jugadores, todos siguen una misma l\'ogica general. Cada script se encarga de consultar uno o varios \textit{endpoints} de la API, parametrizando la petici\'on seg\'un la temporada, el identificador del partido, el equipo o el jugador.

La t\'ecnica empleada se basa en peticiones HTTP mediante la librer\'ia \texttt{requests}. Para cada consulta se construye din\'amicamente la URL del \textit{endpoint}, se incorporan las cabeceras de autenticaci\'on requeridas por la API y se env\'ia la petici\'on. La respuesta recibida se transforma a formato JSON y se almacena en disco, manteniendo una copia de los datos originales antes de aplicar cualquier transformaci\'on.

Un esquema general del proceso seguido por los scripts de extracci\'on es el siguiente:

\begin{enumerate}
    \item Definir las temporadas o identificadores que se desean extraer.
    \item Construir el \textit{endpoint} correspondiente de la API.
    \item Enviar la petici\'on HTTP con las cabeceras de autenticaci\'on.
    \item Recibir la respuesta en formato JSON.
    \item Validar si la respuesta contiene datos.
    \item Guardar el JSON en la carpeta correspondiente.
\end{enumerate}

Esta organizaci\'on permite separar claramente la extracci\'on de la transformaci\'on. Los datos brutos se conservan en el directorio \texttt{DATOS/}, agrupados seg\'un el tipo de informaci\'on y la temporada correspondiente. Por ejemplo, las clasificaciones se almacenan en \texttt{DATOS/clasificaciones/}, los partidos en \texttt{DATOS/partidos\_base/}, las estad\'isticas en \texttt{DATOS/partidos\_stats/} y los eventos en \texttt{DATOS/partidos\_eventos/}.

Adem\'as, los ficheros se organizan por temporada, lo que permite trabajar de forma incremental. Si se desea incorporar una nueva temporada, basta con ejecutar los scripts de extracci\'on para ese a\~no concreto y almacenar los nuevos JSON en la carpeta correspondiente. Esta estructura facilita la actualizaci\'on del sistema y permite mantener un hist\'orico completo de las temporadas analizadas.

\subsection{Tranformación}

La fase de transformación se ha desarrollado principalmente en el directorio \path{ETL/}. Su objetivo es convertir los datos brutos extraídos desde la API, almacenados inicialmente en formato JSON dentro de \path{DATOS/}, en un conjunto de tablas estructuradas y preparadas para su posterior carga en base de datos.

En esta etapa se produce el cambio más importante del proceso ETL: los datos dejan de tener una estructura anidada propia de la API, en formato JSON, y pasan a organizarse en un esquema basado en dimensiones y hechos. Para ello, los scripts Python leen los JSON descargados, recorren sus estructuras internas, seleccionan los campos relevantes, normalizan nombres, convierten tipos de datos, gestionan valores nulos y finalmente generan ficheros CSV dentro del directorio \path{ETL/DSA/}. El resultado de esta fase es una capa analítica intermedia formada por tablas limpias, homogéneas y listas para cargarse en PostgreSQL.

Cada script ha recibido un tratamiento distinto en función de los datos del JSON que debía transformar, aunque todos ellos siguen una serie de pasos comunes con el objetivo de mantener una estructura uniforme en el proceso. Estos pasos son los siguientes:

\begin{enumerate}
    \item Lectura de los ficheros JSON generados en la fase de extracción.
    \item Recorrido de estructuras anidadas de la API.
    \item Selección de campos relevantes para el modelo BI.
    \item Conversión de los datos a un \texttt{DataFrame} de pandas.
    \item Limpieza de valores nulos o inconsistentes.
    \item Conversión de tipos de datos numéricos, fechas e identificadores.
    \item Generación de claves de relación entre tablas.
    \item Exportación final a CSV con separador ;.
\end{enumerate}

Una vez explicado el procedimiento general, se describen las transformaciones realizadas sobre los distintos JSON para construir las tablas que conforman el esquema de la base de datos.

Primero se describen las tablas de dimensiones, las cuales contienen información descriptiva de las entidades principales del sistema, como equipos, jugadores, partidos y tiempo.

La tabla \texttt{dim\_equipos.csv} se genera a partir de los datos básicos de equipos. Esta transformación es sencilla, ya que principalmente se extraen atributos descriptivos del JSON, como el identificador del equipo, nombre, código, país, estadio, ciudad o capacidad. El cambio principal consiste en pasar de una estructura anidada a una fila tabular por equipo.

La tabla \texttt{dim\_jugadores.csv} sigue una lógica similar. Se extraen datos personales y físicos del jugador, como nombre, edad, fecha de nacimiento, nacionalidad, altura, peso y fotografía. En este caso, la transformación también permite unificar la información de jugadores que pueden aparecer en distintas temporadas.

La tabla \texttt{dim\_partidos.csv} requiere una transformación más relevante, ya que se separa la información del partido, los equipos participantes y el resultado. Además, se distingue entre partidos completados e incompletos, evitando perder partidos futuros o sin resultado definitivo. Durante esta transformación, la jornada, que llega desde la API como texto, se convierte en un valor numérico. También se convierten los goles a valores enteros y se normaliza el formato de fecha para que pueda ser utilizado correctamente en la base de datos.

La tabla \texttt{dim\_tiempo.csv} se genera a partir de las fechas de los partidos. En este caso se crea una dimensión temporal que permite analizar la información por año, mes, día y jornada. Esta tabla facilita la realización de consultas temporales y permite agrupar los datos deportivos según diferentes niveles de detalle.

\subsubsection{Transformaciones sobre tablas de hechos}

Las tablas de hechos almacenan medidas deportivas asociadas a partidos, jugadores, equipos o jornadas.

La tabla \texttt{h\_jugador\_partido.csv} es una de las transformaciones más importantes. La API devuelve las estadísticas de jugadores dentro de estructuras anidadas por partido, equipo, jugador y bloque estadístico. El script transforma esta información en una fila por jugador y partido, recogiendo datos como el identificador del partido, el jugador, el equipo, la posición, los minutos disputados, la nota, los goles, las asistencias, los pases y los duelos.

Además, se normalizan posiciones y se convierten las estadísticas a valores numéricos. Los valores ausentes se sustituyen por cero cuando representan acciones que no se han producido durante el partido, como goles, asistencias, disparos o duelos ganados. De esta forma se garantiza que los datos puedan ser utilizados posteriormente en cálculos estadísticos y modelos de minería de datos.

La tabla \texttt{h\_jugador\_temporada.csv} recoge estadísticas acumuladas del jugador durante una temporada. En esta transformación se traducen las posiciones originales de la API al vocabulario usado en el sistema, permitiendo trabajar con categorías más comprensibles para el usuario, como portero, defensa, mediocentro o delantero. También se convierten las métricas acumuladas a tipos numéricos para facilitar su explotación analítica.

La tabla \texttt{h\_equipo\_partido.csv} transforma las estadísticas de equipo por partido. En este caso, una transformación importante es el mapeo de nombres originales de la API a nombres de columnas más claros y homogéneos. Por ejemplo, estadísticas como tiros a puerta, posesión, pases totales, porcentaje de pases acertados o goles esperados se adaptan al formato utilizado en la base de datos. También se eliminan caracteres no numéricos, como el símbolo de porcentaje, y se convierten las columnas a enteros o decimales según corresponda.

La tabla \texttt{h\_equipo\_temporada.csv} se genera a partir de las clasificaciones. Su transformación consiste en separar la información general del equipo, sus puntos, posición, diferencia de goles, forma reciente y estadísticas como local y visitante. Es una transformación principalmente de aplanado de JSON a tabla, en la que se pasa de una estructura jerárquica a un formato tabular preparado para consultas analíticas.

La tabla \texttt{h\_equipo\_jornada.csv} es una de las transformaciones más elaboradas, ya que reconstruye la clasificación jornada a jornada a partir de los partidos completados. Para ello se acumulan puntos, victorias, empates, derrotas, goles a favor y goles en contra tras cada jornada. Posteriormente, los equipos se ordenan por puntos, diferencia de goles y goles a favor para obtener la clasificación correspondiente a cada jornada.

La tabla \texttt{h\_partidos\_eventos.csv} transforma los eventos del partido, como goles, tarjetas, sustituciones o acciones revisadas por VAR. La principal transformación consiste en traducir eventos y detalles al castellano, normalizar los minutos e identificar el jugador, el equipo y, cuando corresponde, el asistente o jugador sustituido. De esta forma, los eventos quedan almacenados de manera homogénea y pueden consultarse posteriormente desde la aplicación.

Como resultado, la carpeta \path{.../ETL/DSA/} almacena los CSV resultantes de todas las transformaciones, con la información que posteriormente se cargará en la base de datos.

\subsection{Carga}

La última fase del proceso ETL corresponde a la carga de los datos transformados en la base de datos. Una vez generados los ficheros CSV en el directorio \texttt{ETL/DSA/}, estos se importan en PostgreSQL mediante la instrucción \texttt{COPY FROM}, ejecutada desde terminal.

La elección de \texttt{COPY FROM} se debe a que permite realizar cargas masivas de datos de forma más eficiente que una inserción fila a fila mediante sentencias \texttt{INSERT}. Como el proceso de transformación ya ha generado ficheros tabulares limpios, con columnas definidas y separador constante, la carga mediante \texttt{COPY} resulta adecuada para importar rápidamente las tablas de dimensiones y hechos.

El flujo general de carga es el siguiente:

\begin{enumerate}
    \item Generar los ficheros CSV mediante los scripts del directorio \texttt{ETL/}.
    \item Comprobar que las columnas del CSV coinciden con la estructura de la tabla destino.
    \item Ejecutar la instrucción \texttt{COPY FROM} desde terminal o desde la consola de PostgreSQL.
    \item Cargar primero las tablas de dimensiones.
    \item Cargar después las tablas de hechos, que dependen de las dimensiones.
    \item Verificar que los registros se han insertado correctamente.
\end{enumerate}

El uso de \texttt{COPY FROM} permite mantener una separación clara entre la fase de transformación y la fase de persistencia. Los scripts Python se encargan de limpiar y estructurar los datos, mientras que PostgreSQL se encarga de importar los CSV ya preparados.

El orden de carga es importante debido a las relaciones logicas entre las tablas. En primer lugar se cargan las dimensiones, ya que contienen las entidades principales del modelo. Posteriormente se cargan las tablas de hechos, que hacen referencia a dichas entidades mediante identificadores.

\begin{enumerate}
    \item \texttt{dim\_equipos}
    \item \texttt{dim\_jugadores}
    \item \texttt{dim\_partidos}
    \item \texttt{dim\_tiempo}
    \item \texttt{h\_equipo\_partido}
    \item \texttt{h\_equipo\_temporada}
    \item \texttt{h\_equipo\_jornada}
    \item \texttt{h\_jugador\_partido}
    \item \texttt{h\_jugador\_temporada}
    \item \texttt{h\_partidos\_eventos}
\end{enumerate}

\section{Aplicación móvil}

El código fuente del sistema se ha estructurado siguiendo una arquitectura modular, separando claramente las responsabilidades entre la capa cliente (\textit{frontend}) y la capa servidor (\textit{backend}). Esta organización facilita el mantenimiento, la escalabilidad y la evolución futura del proyecto.

A nivel general, el repositorio principal se compone de dos bloques funcionales diferenciados:

\begin{itemize}
    \item {Frontend:} Aplicación móvil desarrollada con React Native y Expo.
    \item {Backend:} API REST desarrollada con Node.js y Express encargada de la lógica de negocio y acceso a datos.
\end{itemize}

\subsection{Frontend}
\subsubsection{Organización del frontend}

La aplicación móvil se ha estructurado siguiendo un enfoque basado en componentes reutilizables y separación por responsabilidades funcionales.

\begin{lstlisting}[caption={Estructura del frontend}]
frontend/
|
|-- App.js
`-- src/
    |
    |-- assets/
    |-- components/
    |-- context/
    |-- screens/
    `--theme/
\end{lstlisting}

\begin{itemize}
    \item {assets/:} Recursos estáticos de la aplicación, como iconos, logotipos e imágenes predeterminadas.
    
    \item {components/:} Componentes reutilizables de interfaz. Incluye elementos comunes (cabecera, menús, botones) y componentes especializados para jugadores y partidos.
    
    \item {context/:} Implementación del sistema global de estado mediante React Context. En este caso, se utiliza para la gestión de favoritos del usuario.
    
    \item {screens/:} Pantallas principales de la aplicación, representando cada vista funcional accesible por el usuario.

    \item {theme/:} Código usado para la implementación del modo oscuro.
    
\end{itemize}

\subsubsection{Punto de entrada y navegación}

El punto de entrada principal del frontend se encuentra en el archivo App.js.
Este archivo se encarga de inicializar la navegación de la aplicación, comprobar si existe una sesión guardada, decidir si el usuario debe acceder al login o a la aplicación principal, y envolver la app con los proveedores necesarios.

Para comprobar si existe una sesión guardada, hacemos uso de \texttt{AsyncStorage} con él comprobamos si existe una SESSION\_KEY, si existe, entra a la app directamente, de lo contrario, llevará al usuario a la pantalla de Inicio de Sesión.

Cuando el usuario inicia sesión correctamente y recibe una respuesta correcta del backend, se almacena el objeto usuario en AsyncStorage. Después se actualiza el contexto de favoritos y se reinicia la navegación para llevar al usuario a la aplicación principal. Y para el cierre de sesión simplemente eliminamos la información guaradada en el AsyncStorage y redirigimos al usuario al Login.

En cuanto a la navegación, hemos hecho uso de React Navigation, la cual es la librería más popular y utilizada para gestionar el enrutamiento y la navegación en aplicaciones móviles desarrolladas con React Native. Hemos usado dos tipos de navegadores:
\begin{itemize}
    \item Stack Navigator: para el flujo de las pantallas de Login, Registro y la aplicación principal.

    Este separa claramente las pantallas de acceso de la aplicación auténtica dejando Login y Registro fuera del menú lateral. Una vez que el usuario inicia sesión correctamente, se accede a MainApp, que contiene el Drawer Navigator con las pantallas principales.
    
    \item Drawer Navigator: contiene el menú lateral con las secciones principales de la app. Agrupa las pantallas principales y las pantallas de detalle. Aunque algunas pantallas, como DetalleEquipo o DetalleJugador, no aparecen como opciones visibles del menú, se registran en el navegador para poder acceder a ellas desde listados, buscadores o tarjetas.

\end{itemize}

En conjunto, App.js centraliza la navegación principal del frontend. Este archivo controla el acceso inicial según la sesión del usuario, registra las pantallas disponibles, configura el menú lateral y monta elementos globales como el contexto de favoritos y el asistente Agente de Oro.


\subsubsection{Estructura de pantallas}


El frontend organiza las vistas principales de la aplicación dentro de la carpeta screens. Cada archivo representa una pantalla completa de la aplicación y se encarga de coordinar la carga de datos, la interacción del usuario y la navegación hacia otras vistas.

\begin{lstlisting}[caption={Pantallas de la aplicación}]
screens/
|
|-- LoginScreen.js
|-- RegistroScreen.js
|-- InicioScreen.js
|-- TemporadaScreen.js
|-- EquiposScreen.js
|-- JugadoresScreen.js
|-- PartidosScreen.js
|-- DetalleTemporadaScreen.js
|-- DetalleEquipoScreen.js
|-- DetalleJugadorScreen.js
|-- DetallePartidoScreen.js
|-- PerfilScreen.js
`-- AjustesScreen.js
\end{lstlisting}
 Entre ellas se pueden distinguir:
 \begin{itemize}
     \item Pantallas de inicio: Login, Registro e Inicio.
     \item Pantallas principales: Temporada, Equipos, Jugadores y Partidos.
     \item Pantallas de detalles: pantallas en las que entra el análisis de cada pantalla del punto anterior.
     \item Pantalla de configuración del usuario: Perfil y Ajustes.
 \end{itemize}

Ademas de la carpeta screens, la aplicación utiliza la carpeta components para dividir el contenido interno de muchas pantallas. Esto es especialmente importante en las pantallas de detalle, ya que estas no contienen toda la lógica visual en un único archivo, sino que actúan como contenedores de navegación por pestañas.

Dentro de components existen subcarpetas organizadas por dominio funcional: temporadas, equipos, jugadores y partido. Cada una contiene los componentes que se muestran dentro de las pestañas superiores de su pantalla de detalle correspondiente, como las pantallas de análisis, gráficos, plantillas, etc.
\begin{lstlisting}[caption=Pantallas de detalle]
components/
|-- temporadas/
|-- equipos/
|-- jugadores/
`-- partido/
\end{lstlisting}
\subsubsection{Uso de pestañas}
En las pantallas de detalle, la aplicación utiliza pestañas superiores para organizar la información de forma clara. Esta solución permite dividir el contenido en secciones independientes, evitando que toda la información de una temporada, un equipo, un jugador o un partido aparezca concentrada en una única pantalla. Para ello hemos utilizado la libreria \texttt{MaterialTopTabNavigator}.

Cada Tab.Screen representa una pestaña diferente dentro de la pantalla de detalle. En lugar de incluir todo el contenido directamente en la pantalla principal, cada pestaña carga un componente independiente. Por ejemplo, en el detalle de un equipo, las pestañas muestran información general, partidos, trayectoria, clasificación, plantilla y análisis.

\begin{table}[H]
\centering
\caption{Pestañas principales de las pantallas de detalle}
\label{tab:pestanas_detalle}
\renewcommand{\arraystretch}{1.2}
\begin{tabular}{|p{5cm}|p{10cm}|}
\hline
\textbf{Pantalla de detalle} & \textbf{Pestañas principales} \\ \hline

\texttt{DetalleTemporadaScreen} &
Partidos, Clasificación, Información, Rankings, Análisis y Equipos. \\ \hline

\texttt{DetalleEquipoScreen} &
Información, Partidos, Trayectoria, Clasificación, Plantilla y Análisis. \\ \hline

\texttt{DetalleJugadorScreen} &
Información, Partidos, Trayectoria, Estadísticas y Atributos. \\ \hline

\texttt{DetallePartidoScreen} &
Previa, Alineación, Eventos y Post-partido. \\ \hline

\end{tabular}
\end{table}

El uso de pestañas mejora la organización de la interfaz y hace que la navegación
sea mas cómoda. Cada pantalla de detalle actúa como contenedor, mientras que los
componentes internos de cada pestaña se encargan de mostrar y gestionar una parte
concreta de la información.
\subsubsection{Conexión con backend}
La comunicación entre el frontend y el backend se realiza mediante peticiones
HTTP. El frontend, desarrollado con React Native y Expo, utiliza la función
fetch para solicitar datos a la API REST del backend, que responde en formato
JSON. Todas las peticiones se construyen a partir de la variable
EXPO\_PUBLIC\_API\_URL. De esta forma, la aplicación puede conectarse al backend
local durante el desarrollo, o a otro servidor en caso de despliegue, cambiando
unicamente la configuración del entorno.

El flujo de ejecución, comienza con el usuario interactuando con la pantalla, como por ejemplo, ver los datos de una temporada en concreto en su pantalla correspondiente. Al pulsar el botón el frontend realiza una petición \textit{fetch} al backend.

Este tipo de petición permite solicitar al backend datos concretos, por ejemplo,
los partidos de una temporada y jornada determinadas. El frontend construye la
URL con los parámetros seleccionados por el usuario, y el backend los interpreta
para realizar la consulta correspondiente.

Una vez realizada la petición, se llama al controlador para que se ejecute el endpoint, consulte en la base de datos y devuelva el resultado mediante un JSON, asi poder mostrar la información al usuario con el diseño que hayamos implementado.

Estas peticiones se validan mediante estados de carga que permiten que la interfaz responda correctamente mientras espera la respuesta del servidor. Si la petición se completa con éxito,
los datos se almacenan en el estado de la pantalla. Si ocurre un error, se muestra
un mensaje informativo al usuario.
\subsubsection{Componentes reutilizables}
El frontend no esta formado unicamente por pantallas completas, sino también por
componentes reutilizables. Estos componentes se encuentran dentro de la carpeta
components, y permiten evitar duplicar código en distintas partes de la aplicación. Entre estos componentes encontramos:
\begin{itemize}
    \item Menú lateral: El componente menu.js implementa el contenido del menú lateral de la aplicación. Desde este menú, el usuario puede acceder a las secciones principales, como Inicio,
Temporadas, Equipos, Jugadores y Partidos. También incluye accesos al perfil,
ajustes y cierre de sesión.
    \item Botón de favorito: FavoritoButton.js es un componente reutilizable que permite marcar o desmarcar equipos y jugadores como favoritos. Se utiliza tanto en listados como en pantallas
de detalle, y muestra un icono de corazon que cambia segun el estado del favorito.
    \item Header: El componente header.js define la cabecera superior utilizada en las pantallas
principales. Incluye el botón para abrir el menú lateral, el titulo de la pantalla
y el acceso al buscador global. Asi todas las pantallas tienen el mismo header sin tener que repetir código.
    \item Agente de oro (Chatbot): El componente AgenteDeOro.js representa el asistente flotante de la aplicación. Se monta de forma global desde App.js, y su visibilidad depende de la pantalla en
la que se encuentra el usuario. Aunque visualmente aparece como un elemento independiente, se considera un componente reutilizable porque no pertenece a una pantalla concreta, sino que
puede mostrarse en varias secciones de la aplicación.
\end{itemize}
\subsubsection{Visualización de datos}
Una parte importante del frontend es la representación visual de estadísticas.
La aplicación no se limita a mostrar listados de datos, sino que también incluye
gráficos, rankings, tablas y tarjetas resumen, con el objetivo de facilitar el
análisis de temporadas, equipos, jugadores y partidos. La idea de esto, es de crear dashboards para cada módulo de la aplicación; temporada, equipo, jugador y partido. De forma que el usuario pueda ver de forma analítica y visual estadísticas importantes en cada aspecto, y que pueda interactuar con ellos filtrando por equipo, edad y demás atributos.

Las librerias que nos han permitido construir estos gráficos directamente dentro de la aplicación son las siguientes:
\begin{itemize}
    \item victory-native \cite{victorynative}: Gráficos de barras, lineas, dispersión y radar.
    \item react-native-gifted-charts \cite{giftedcharts}: Gráficos específicos, como barras y radar de atributos.
    \item react-native-svg \cite{reactnativesvg}: Soporte gráfico para renderizar visualizaciones.
\end{itemize}

La construcción de los gráficos sigue un patrón común en la aplicación. En primer lugar, el frontend obtiene los datos desde el backend mediante una petición HTTP. Después, estos datos se almacenan en el estado del componente y se transforman al formato que necesita la librería de gráficos. Finalmente, se renderiza el gráfico mediante componentes de Victory, como \texttt{VictoryChart}, \texttt{VictoryAxis} o
\texttt{VictoryScatter}.

El componente \texttt{VictoryChart} actúa como contenedor principal del gráfico. Dentro de el se definen los ejes mediante \texttt{VictoryAxis}. El primer eje representa el eje horizontal, mientras que el eje marcado como
\texttt{dependentAxis} representa el eje vertical. Finalmente,
\texttt{VictoryScatter} recibe los datos que se quieren representar y dibuja un punto por cada elemento del array.

Para que un gráfico pueda representarse correctamente, los datos deben adaptarse previamente al formato esperado por la librería. En los gráficos de dispersión, por ejemplo, cada elemento del array debe incluir una coordenada \texttt{x} y una coordenada \texttt{y}. Estas coordenadas se calculan a partir de los campos recibidos desde el backend.

Este patrón se repite en los distintos dashboards de la aplicación. Primero, se obtienen los datos desde la API; después, se limpian o convierten los valores numéricos; a continuación, se construye un array adaptado al tipo de gráfico; y, por último, se pasa dicho array al componente visual correspondiente.
\begin{itemize}
\item{Dashboard de temporada:} 
En el módulo de temporadas, los gráficos permiten analizar el rendimiento general de equipos y jugadores a lo largo de una temporada. Para este he querido reflejar el rendimiento en general, defensa, creación y ataque, ya que son los aspectos mas importantes y relevante de un equipo.

En la tabla ~\ref{tab:graficos_dashboard_temporada} podemos ver definidos los diferentes gráficos que se han decidido implementar y que pueden dar una buena visión de como ha sido la temporada tanto en general, como para cada equipo, ya que el dashboard incluye un filtro para poder ver los datos de un equipo en específico.

\begin{longtable}{|p{5cm}|p{9cm}|}
\caption{Justificación de los principales gráficos del dashboard de temporada}
\label{tab:graficos_dashboard_temporada} \\

\hline
\textbf{Gráfico / bloque} & \textbf{Justificación} \\ \hline
\endfirsthead

\hline
\textbf{Gráfico / bloque} & \textbf{Justificación} \\ \hline
\endhead

\hline
\multicolumn{2}{r}{\textit{Continúa en la siguiente página}} \\ \hline
\endfoot

\hline
\endlastfoot

KPIs globales &
Resume la temporada mediante indicadores clave que permiten conocer rápidamente el estado general de la competición. \\ \hline

Rendimiento pasado &
Compara el rendimiento de los equipos según su diferencia de goles, identificando los mejores y peores registros. \\ \hline

Cuadrante de equipos &
Relaciona los goles a favor y los goles en contra para visualizar el equilibrio ofensivo y defensivo de cada equipo. \\ \hline

Rendimiento de jugadores &
Compara los minutos disputados con la nota media obtenida por cada jugador durante la temporada. \\ \hline

Acciones defensivas &
Representa métricas defensivas como bloqueos, entradas e intercepciones para evaluar el rendimiento sin balón. \\ \hline

Líderes del mediocampo &
Destaca a los jugadores más relevantes en la creación de juego mediante estadísticas de pases y asistencias. \\ \hline

Delanteros &
Relaciona el número de tiros realizados con los goles anotados para analizar la eficacia ofensiva de los atacantes. \\ \hline
\end{longtable}

\begin{figure}[H]
\centering
  \begin{subfigure}[b]{0.4\textwidth} % 0.4 es el tamaño de la imagen
    \includegraphics[width=0.9\textwidth]{figures/dash_tempo1.jpeg} % 0.9 es el tamaño dentro del espacio anterior
    \caption{KPI's Iniciales} % Título de la imagen
    \label{fig:dt1} % Etiqueta de la imagen
  \end{subfigure}
  %
  \begin{subfigure}[b]{0.4\textwidth} %Para incluir una segunda figura
    \includegraphics[width=1\textwidth]{figures/dash_tempo2.jpeg}
    \caption{Cuadrante de goles y Rendimiento/minutos}
    \label{fig:dt2}
  \end{subfigure}
    \caption{Gráficos de temporada 1}
  \label{fig:dash_tempo1}
  %
\end{figure}
\begin{figure}[H]
\centering
\begin{subfigure}[b]{0.4\textwidth} % 0.4 es el tamaño de la imagen
    \includegraphics[width=1\textwidth]{figures/dash_tempo3.jpeg} % 0.9 es el tamaño dentro del espacio anterior
    \caption{Rendimiento defensivo y creativo} % Título de la imagen
    \label{fig:dt3} % Etiqueta de la imagen
  \end{subfigure}
  %
  \begin{subfigure}[b]{0.4\textwidth} %Para incluir una segunda figura
    \includegraphics[width=1\textwidth]{figures/dash_tempo4.jpeg}
    \caption{Rendimiento ofensivo}
    \label{fig:dt4}
  \end{subfigure}
  \caption{Gráficos de temporada 2}
  \label{fig:dash_tempo2}
\end{figure}

\item {Dashboard de partido:}
El dashboard de partido permite interpretar lo ocurrido en un encuentro desde
varios puntos de vista. Primero, se ofrece una comparativa general entre equipos, para entender el equilibrio del partido. Después, los gráficos individuales permiten analizar que jugadores fueron mas eficaces en ataque, cuales generaron mas juego y quienes tuvieron mayor aportación defensiva. Esta división facilita pasar de una visión colectiva del partido a una lectura individual del rendimiento.

En la Tabla~\ref{tab:graficos_dashboard_partido} podemos ver las definiciones de los gráficos implementadas para analizar adecuadamente cada partido:

\begin{longtable}{|p{5cm}|p{9cm}|}
\caption{Justificación de los principales gráficos del dashboard de partido}
\label{tab:graficos_dashboard_partido} \\

\hline
\textbf{Gráfico / bloque} & \textbf{Justificación} \\ \hline
\endfirsthead

\hline
\textbf{Gráfico / bloque} & \textbf{Justificación} \\ \hline
\endhead

\hline
\multicolumn{2}{r}{\textit{Continúa en la siguiente página}} \\ \hline
\endfoot

\hline
\endlastfoot

Comparativa de equipos &
Resume las principales estadísticas del encuentro para obtener una visión general del rendimiento de ambos equipos. \\ \hline

Radar de dominio &
Compara el grado de dominio del partido entre los dos equipos mediante indicadores de posesión, pases y acciones ofensivas. \\ \hline

Finalización &
Analiza la capacidad ofensiva de los equipos a través de los tiros realizados y su precisión. \\ \hline

Eficacia ofensiva &
Relaciona el número de tiros con los goles anotados para evaluar la efectividad de los jugadores en ataque. \\ \hline

Creadores &
Destaca a los jugadores con mayor capacidad de generación de juego mediante pases clave y asistencias. \\ \hline

Muralla defensiva &
Resume el rendimiento defensivo mediante estadísticas de entradas, bloqueos e intercepciones. \\ \hline

Duelos y regates &
Representa el rendimiento individual en acciones de uno contra uno, tanto ofensivas como defensivas. \\ \hline

Notas por posición &
Compara las valoraciones obtenidas por los jugadores en función de su posición táctica sobre el terreno de juego. \\ \hline

Top impacto &
Ordena los jugadores con mayor influencia en el desarrollo del partido según su rendimiento global. \\ \hline

Porteros &
Evalúa la actuación de los guardametas mediante estadísticas de paradas y goles encajados. \\ \hline

\end{longtable}


\begin{figure}[H]
\centering
  \begin{subfigure}[b]{0.4\textwidth} % 0.4 es el tamaño de la imagen
    \includegraphics[width=0.9\textwidth]{figures/dash_partido1.jpeg} % 0.9 es el tamaño dentro del espacio anterior
    \caption{Comparativa de equipos} % Título de la imagen
    \label{fig:dp1} % Etiqueta de la imagen
  \end{subfigure}
  %
  \begin{subfigure}[b]{0.4\textwidth} %Para incluir una segunda figura
    \includegraphics[width=1\textwidth]{figures/dash_partido2.jpeg}
    \caption{Radar de dominio}
    \label{fig:dp2}
  \end{subfigure}
  %
\begin{subfigure}[b]{0.4\textwidth} % 0.4 es el tamaño de la imagen
    \includegraphics[width=1\textwidth]{figures/dash_partido3.jpeg} % 0.9 es el tamaño dentro del espacio anterior
    \caption{Eficacia ofensiva} % Título de la imagen
    \label{fig:dp3} % Etiqueta de la imagen
  \end{subfigure}
  %
  \begin{subfigure}[b]{0.4\textwidth} %Para incluir una segunda figura
    \includegraphics[width=1\textwidth]{figures/dash_partido4.jpeg}
    \caption{Top impacto}
    \label{fig:dp4}
  \end{subfigure}
  \caption{Gráficos de partido}
  \label{fig:dash_tempo}
\end{figure}
\item Dashboard de equipo:

El dashboard de equipo está orientado al análisis del rendimiento de un club durante una temporada. Para facilitar la interpretación de la información, el análisis se divide en dos aspectos principales:

\begin{itemize}
    \item Análisis colectivo: estudia el comportamiento del equipo como conjunto. En este apartado se incluyen indicadores generales del rendimiento, la evolución de la posición durante la temporada, el perfil de juego representado mediante un gráfico radar, el rendimiento frente a distintos tipos de rivales, la comparativa entre partidos como local y visitante, la dinámica de los últimos cinco encuentros y las recomendaciones de fichajes obtenidas mediante los modelos de minería de datos.

    \item Análisis individual: estará orientado al estudio del rendimiento de los jugadores que forman parte de la plantilla. Su objetivo será ofrecer una visión detallada de las estadísticas individuales de cada futbolista, permitiendo identificar fortalezas, debilidades y su aportación al rendimiento global del equipo.
\end{itemize}

La Tabla~\ref{tab:dashboard_equipo} resume la finalidad de cada uno de los gráficos que componen el análisis colectivo del dashboard de equipo.

\begin{table}[H]
\centering
\caption{Justificación de los gráficos del dashboard de equipo}
\label{tab:dashboard_equipo}
\renewcommand{\arraystretch}{1.2}
\begin{tabular}{|p{6cm}|p{9cm}|}
\hline
\textbf{Gráfico / bloque} & \textbf{Justificación} \\ \hline

KPIs del equipo &
Resume la información principal del equipo, incluyendo posición, puntos y estadísticas generales de la temporada. \\ \hline

Posición durante la temporada &
Representa la evolución del equipo en la clasificación jornada a jornada, permitiendo identificar cambios de rendimiento. \\ \hline

Radar del equipo &
Resume el perfil de juego del equipo mediante seis dimensiones que sintetizan sus principales fortalezas y debilidades. \\ \hline

Rendimiento por tipo de rival &
Compara los resultados obtenidos frente a rivales de distintos niveles competitivos para evaluar el comportamiento del equipo en diferentes contextos. \\ \hline

Comparativa local vs. visitante &
Analiza las diferencias de rendimiento cuando el equipo juega en casa o fuera de ella. \\ \hline

Últimos cinco partidos &
Muestra la dinámica reciente del equipo y permite acceder directamente al detalle de cada encuentro. \\ \hline

Recomendación de fichajes &
Presenta los jugadores recomendados para reforzar la plantilla a partir de los modelos de minería de datos desarrollados. \\ \hline

\end{tabular}
\end{table}

\item {Dashboard de jugador:}

\end{itemize}

\subsubsection{Estados e interacción}

Además de mostrar información, el \textit{frontend} debe permitir la interacción del usuario y adaptarse al estado de cada pantalla. Para ello, la aplicación utiliza mecanismos propios de React Native que permiten gestionar la carga de datos, los errores, la navegación, los filtros, los modales y las acciones realizadas por el usuario.

Uno de los aspectos más importantes es la gestión de las peticiones al \textit{backend}. Cuando una pantalla solicita información, la aplicación debe contemplar distintos escenarios: que los datos se estén cargando, que la respuesta sea correcta o que se produzca algún error. En función de esta situación, la interfaz muestra un indicador de carga, un mensaje informativo o el contenido correspondiente. Este patrón se utiliza en pantallas como los listados de equipos y jugadores, los detalles de temporada, los análisis de partido y los \textit{dashboards} estadísticos.

La interacción también se refleja en el uso de listas y tarjetas pulsables. Estas permiten acceder desde una vista general a una pantalla de detalle, enviando internamente el identificador necesario para cargar la información concreta del equipo, jugador, partido o temporada seleccionada. De esta forma, la navegación resulta sencilla para el usuario y mantiene una estructura coherente dentro de la aplicación.

Asimismo, se utilizan modales, buscadores y filtros para facilitar la consulta de información sin sobrecargar la interfaz. Los modales permiten seleccionar opciones como temporadas, jornadas o equipos sin abandonar la pantalla actual, mientras que los campos de búsqueda y filtros ayudan a localizar rápidamente elementos concretos dentro de listas extensas.

Por último, la aplicación incorpora la gestión de favoritos, permitiendo al usuario marcar o desmarcar equipos y jugadores. Esta acción actualiza la interfaz y se sincroniza con el \textit{backend} para conservar la información de forma persistente.

En conjunto, la gestión de estados e interacción permite que la aplicación sea dinámica, clara y reactiva, ofreciendo una experiencia más fluida frente a cambios realizados por el usuario, retrasos del servidor o errores de conexión.

\subsubsection{Simulación manual de una temporada}

El sistema incorpora una pantalla específica para simular manualmente una temporada. Esta funcionalidad permite modificar los resultados de los partidos pendientes, recalcular la clasificación y ejecutar una simulación probabilística sobre el escenario creado.

La pantalla se integra como una sección independiente dentro del menú lateral y se organiza en dos pestañas: \texttt{RESULTADOS} (véase Figura~\ref{fig:resultados}) y \texttt{CLASIFICACIÓN} (véase Figura~\ref{fig:clasificacion}). En la primera, el usuario puede desplazarse entre jornadas y editar únicamente los encuentros que todavía no se han disputado. Los partidos ya completados permanecen bloqueados para evitar cambios en los datos reales.

Los resultados introducidos se conservan temporalmente durante la simulación, incluso al cambiar de jornada. Al pulsar \texttt{Recalcular escenario} (véase Figura~\ref{fig:boton_sim}), el sistema parte de la clasificación real y actualiza los partidos jugados, victorias, empates, derrotas, goles y puntos. Posteriormente, ordena la nueva clasificación según los criterios de puntos, diferencia de goles y goles a favor.

\begin{figure}[H]
\centering
  \begin{subfigure}[b]{0.4\textwidth} % 0.4 es el tamaño de la imagen
    \includegraphics[width=0.9\textwidth]{figures/sim_result.jpeg} % 0.9 es el tamaño dentro del espacio anterior
    \caption{Editor de resultados} % Título de la imagen
    \label{fig:sim_result} % Etiqueta de la imagen
  \end{subfigure}
  %
  \begin{subfigure}[b]{0.4\textwidth} %Para incluir una segunda figura
    \includegraphics[width=1\textwidth]{figures/sim_botonrec.jpeg}
    \caption{Botón \texttt{Recalcular escenarios}}
    \label{fig:boton_sim}
  \end{subfigure}
  \caption{Pantalla \texttt{RESULTADOS}}
  \label{fig:resultados}
\end{figure}

En la pestaña \texttt{CLASIFICACIÓN}, el usuario puede consultar tanto la tabla simulada (véase Figura~\ref{fig:sim_clas}) como las probabilidades calculadas mediante el algoritmo de Monte Carlo (véase Figura ~\ref{fig:sim_prob}). Este proceso se ejecuta de forma temporal sobre el escenario definido, sin modificar la clasificación ni las probabilidades oficiales almacenadas en la base de datos.

\begin{figure}[H]
\centering
  \begin{subfigure}[b]{0.4\textwidth} % 0.4 es el tamaño de la imagen
    \includegraphics[width=0.9\textwidth]{figures/sim_clasi.jpeg} % 0.9 es el tamaño dentro del espacio anterior
    \caption{Pantalla de clasificación} % Título de la imagen
    \label{fig:sim_clas} % Etiqueta de la imagen
  \end{subfigure}
  %
  \begin{subfigure}[b]{0.4\textwidth} %Para incluir una segunda figura
    \includegraphics[width=0.9\textwidth]{figures/sim_prob.jpeg}
    \caption{Pantalla de probabilidades}
    \label{fig:sim_prob}
  \end{subfigure}
  \caption{Pantalla \texttt{PROBABILIDADES}}
  \label{fig:clasificacion}
\end{figure}


Esta funcionalidad aporta una capa interactiva al sistema, ya que permite al usuario construir escenarios alternativos sin alterar los datos reales de la temporada. El simulador permite analizar cómo determinados resultados podrían afectar a la clasificación y a las probabilidades finales de cada equipo, lo que lo convierte en una herramienta útil para explorar posibles desenlaces de la competición.

\subsubsection{Modo oscuro}

La aplicación incorpora un modo oscuro que puede activarse manualmente desde la pantalla de ajustes. El usuario puede alternar entre el tema claro y el oscuro mediante un interruptor, y la preferencia se almacena localmente con \texttt{AsyncStorage}, por lo que se conserva al cerrar la aplicación o la sesión.

La implementación utiliza un sistema global de tematización formado por un \texttt{ThemeProvider} y el hook \texttt{useTheme()}. Este contexto proporciona el estado del tema, la paleta de colores activa y la función para cambiar de modo, permitiendo que todas las pantallas y componentes utilicen una configuración visual común.

El tema se aplica a la navegación, el menú lateral, los formularios, los modales, las tarjetas y otros componentes reutilizables. También se adapta la \texttt{StatusBar} para mantener una apariencia coherente en toda la aplicación.

Las gráficas y tablas de probabilidades requieren una adaptación adicional, ya que deben modificarse los colores de los ejes, etiquetas, leyendas y textos internos para garantizar un contraste adecuado. Las tablas de Monte Carlo mantienen sus colores informativos, pero ajustan el texto para asegurar la legibilidad en ambos temas.

En conjunto, el modo oscuro mejora la accesibilidad visual y la comodidad de uso en entornos con poca iluminación. Además, su gestión centralizada permite aplicarlo de manera uniforme sin modificar la lógica funcional de los componentes.

\subsection{Backend}
% En este apartado explicaremos las diferentes partes que hemos implementado en el backend para que la app y la base de datos pueden comunicarse de forma eficiente y escalable. 

% \subsection{Arranque del servidor}

% El punto de entrada del backend se encuentra en el archivo app.js. Este archivo es el encargado de inicializar el servidor Express, cargar las variables de entorno, configurar los middlewares principales y montar las rutas de la API.

% En primer lugar, se carga la configuración definida en el archivo .env mediante la librería dotenv. Esto permite utilizar variables de entorno como el puerto del servidor o los datos de conexión necesarios para otros módulos del backend. A continuación, se crea una instancia de Express y se configuran los middlewares globales.

% Uno de los middlewares utilizados es cors, que permite que el frontend pueda
% realizar peticiones al backend aunque ambos se ejecuten en orígenes distintos.
% También se utiliza express.json(), necesario para que el servidor pueda recibir y procesar cuerpos de petición en formato JSON, por ejemplo durante el inicio de sesión, el registro de usuarios o las consultas al chatbot.

% Todas las rutas de la API se agrupan bajo el prefijo /api. Esto permite mantener una estructura clara y diferenciar las peticiones dirigidas al backend del resto de recursos de la aplicación.

% \begin{lstlisting}[caption={app.js}]
% require('dotenv').config();

% const express = require('express');
% const cors = require('cors');

% const routes = require('./routes');
% const app = express();

% app.use(cors());
% app.use(express.json());

% app.use('/api', routes);

% const PORT = process.env.PORT || 3000;
% app.listen(PORT, () => {
%   console.log(`Servidor corriendo en el puerto ${PORT}`);
% });
% \end{lstlisting}

% En este fragmento se observa que el servidor utiliza el puerto indicado en la
% variable de entorno PORT. En caso de que no exista dicha variable, se utiliza
% por defecto el puerto 3000. Finalmente, mediante app.listen(), el servidor queda a la espera de recibir peticiones HTTP.

% Por tanto, app.js centraliza el arranque del backend: carga la configuracion,
% prepara Express, habilita la comunicacion con el frontend, registra las rutas
% bajo el prefijo /api y deja el servidor escuchando peticiones.

% \subsection{Conexión con Base de Datos}

% La conexion del backend con la base de datos se centraliza en el archivo
% config/db.js. Para ello se utiliza la libreria pg, que permite conectar una
% aplicación Node.js con una base de datos PostgreSQL.

% En este archivo se crea un objeto Pool. Un pool de conexiones permite reutilizar conexiones a la base de datos en lugar de abrir una nueva conexión en cada petición. Esto mejora el rendimiento del servidor y permite gestionar de forma mas eficiente múltiples consultas simultaneas realizadas desde el frontend.

% \begin{lstlisting}[caption = Configuración de la base de datos]
% const { Pool } = require('pg');

% console.log("Conectando a DB con usuario:", process.env.DB_USER);

% const pool = new Pool({
%   user: String(process.env.DB_USER || ''),
%   host: String(process.env.DB_HOST || ''),
%   database: String(process.env.DB_NAME || ''),
%   password: String(process.env.DB_PASSWORD || ''),
%   port: process.env.DB_PORT,
% });

% module.exports = pool;
% \end{lstlisting}

% Como vemos, recibimos la configuración de la base de datos Postgres desde las variables de entorno para realizar la conexión. Una vez creado el pool, este se exporta mediante module.exports para poder ser utilizado desde los controladores y servicios del backend.

% \subsection{Rutas y controladores}

% Una parte fundamental de la implementación del backend es la separación entre
% rutas y controladores. Las rutas definen los endpoints disponibles de la API,
% mientras que los controladores contienen la lógica que se ejecuta cuando el
% frontend realiza una petición.

% Esta separación permite que el código sea mas claro y mantenible. Por un lado,
% los archivos de routes indican que URL existe y que método HTTP utiliza. Por
% otro lado, los archivos de controllers se encargan de consultar la base de datos, procesar los parámetros recibidos y devolver una respuesta JSON.

% El archivo routes/index.js actua como punto central de las rutas del backend.
% Cada grupo de endpoints se monta bajo un prefijo concreto. Como el archivo app.js monta todas las rutas bajo /api, una ruta definida como /equipos pasa a estar disponible finalmente como /api/equipos.

% \begin{lstlisting}[caption=routes/index.js]
% const express = require('express');
% const router = express.Router();

% const equiposRoutes = require('./equipos');
% const partidosRoutes = require('./partidos');
% const UsuariosRoutes = require('./usuarios');
% const jugadoresRoutes = require('./jugadores');
% const temporadasRoutes = require('./temporadas');
% const busquedaRoutes = require('./buscador');
% const favoritosRoutes = require('./favoritos');
% const chatRoutes = require('./chatbot');

% router.use('/equipos', equiposRoutes);
% router.use('/partidos', partidosRoutes);
% router.use('/usuarios', UsuariosRoutes);
% router.use('/temporadas', temporadasRoutes);
% router.use('/jugadores', jugadoresRoutes);
% router.use('/buscador', busquedaRoutes);
% router.use('/favoritos', favoritosRoutes);
% router.use('/chatbot', chatRoutes);

% module.exports = router;
% \end{lstlisting}

% Podemos ver un claro ejemplo de ruta y controlador en la ruta de /equipos, una de las mas importantes por su gran número de endpoints.

% \begin{lstlisting}[caption=routes/equipos.js]
% const express = require('express');
% const router = express.Router();
% const equipoController = require('../controllers/equipoController');

% router.get('/', equipoController.getEquipos);
% router.get('/:id', equipoController.getEquipoPorId);
% router.get('/info/:id_equipo', equipoController.getInfoEquipo);
% router.get('/partidos/:id_equipo', equipoController.getPartidosEquipoPorId);
% //MAS ENDPOINTS...

% module.exports = router;
% \end{lstlisting}

% \begin{lstlisting}[caption= función de controllers/EquipoController.js]
% const pool = require('../config/db');

% const getEquipos = async (req, res) => {
%   try {
%     const result = await pool.query(
%       'SELECT id_equipo, nombre_equipo, logo FROM dim_equipo ORDER BY nombre_equipo'
%     );

%     res.json(result.rows);
%   } catch (error) {
%     console.error('Error al obtener los equipos:', error);
%     res.status(500).json({ error: 'Error interno del servidor' });
%   }
% };
% \end{lstlisting}

% El flujo de trabajo sería:
% \begin{enumerate}
%     \item Frontend hace petición a endpoint, por ejemplo, /api/equipos
%     \item Eso lleva a index.js, que redirige la llamada a equipos.js
%     \item equipos.js llama a la función asociada a la ruta, en este caso getEquipos
%     \item Consulta BD PostgreSQL y devuelve JSON con la información.
% \end{enumerate}

\subsubsection{Organización del backend}

El backend sigue una arquitectura basada en el patrón {MVC simplificado} (Modelo-Controlador-Rutas), separando la lógica de acceso, procesamiento y exposición de servicios.

\begin{lstlisting}[caption={Estructura del backend}]
backend/src/
|
|-- app.js
|-- config/
|-- controllers/
|-- routes/
|-- services/
|-- utils/
`-- prompts/
\end{lstlisting}

\begin{itemize}
    \item {app.js:} Punto de entrada principal del servidor Express y configuración general de middleware y rutas.
    
    \item {config/:} Módulo encargado de la conexión con PostgreSQL y configuración del acceso a base de datos y servicios externos como la API de Gemini.
    
    \item {controllers/:} Contiene la lógica de negocio asociada a cada recurso del sistema. Cada controlador procesa peticiones HTTP, consulta la base de datos y devuelve respuestas estructuradas.
    
    \item {routes/:} Define los distintos endpoints de la API REST y enlaza cada ruta con su controlador correspondiente.
    
    \item {services/:} Agrupa lógica auxiliar más específica, como generación de SQL o procesamiento de respuestas del \textit{chatbot}.

    \item {utils/:} Contiene funciones de apoyo para el \textit{chatbot}, por ejemplo validación de SQL.
    
    \item {prompts/:} Contiene el \textit{prompt} usado para guiar al \textit{chatbot} en la generación de consultas sobre la base de datos.

    

\end{itemize}

Las carpetas routes y controllers son dos de las partes más importantes del \textit{backend}, ya que permiten separar la definición de los \textit{endpoints} de la lógica que se ejecuta cuando dichos \textit{endpoints} son llamados.

La carpeta routes contiene los archivos que definen las rutas de la API. Es decir,establece que URLs existen, que método HTTP utilizan y que controlador debe ejecutarse en cada caso. Por ejemplo, una ruta puede indicar que una petición GET a /equipos debe ejecutar la función getEquipos del controlador de equipos.

La carpeta controllers contiene la lógica asociada a cada ruta. En estos archivos se reciben los parámetros enviados por el \textit{frontend}, se realizan las consultas a la base de datos, se controlan posibles errores y se devuelve la respuesta en formato JSON.

Esta separación es importante porque mejora la organización del código, facilita el mantenimiento y permite ampliar la API de forma mas ordenada. Si en el futuro se añade una nueva funcionalidad, basta con crear o modificar su ruta correspondiente y asociarla a una función del controlador, sin mezclar la declaración de \textit{endpoints} con la lógica interna de la aplicación.

En la tabla~\ref{tab:endpoints_api} mostramos un resumen de los diferentes \textit{endpoints} que se han implementado en cada ruta:

\renewcommand{\arraystretch}{1.2}

\begin{longtable}{|p{2cm}|p{6.8cm}|p{7cm}|}
\caption{Endpoints implementados en la API REST.}
\label{tab:endpoints_api} \\

\hline
\textbf{Método} & \textbf{Endpoint} & \textbf{Descripción} \\ \hline
\endfirsthead

\hline
\textbf{Método} & \textbf{Endpoint} & \textbf{Descripción} \\ \hline
\endhead

\hline
\multicolumn{3}{r}{\textit{Continúa en la siguiente página}} \\ \hline
\endfoot

\hline
\endlastfoot

GET & \path{/api/equipos} & Obtiene el listado de equipos. \\ \hline
GET & \path{/api/equipos/:id} & Obtiene la información de un equipo concreto. \\ \hline
GET & \path{/api/equipos/info/:id_equipo} & Obtiene el resumen detallado de un equipo. \\ \hline
GET & \path{/api/equipos/partidos/:id_equipo} & Obtiene los partidos de un equipo. \\ \hline
GET & \path{/api/equipos/plantilla/:id_equipo} & Obtiene la plantilla de un equipo por temporada. \\ \hline
GET & \path{/api/equipos/trayectoria/:id_equipo} & Obtiene la trayectoria histórica de un equipo. \\ \hline
GET & \path{/api/equipos/dashboard/:id_equipo} & Obtiene los datos agregados del dashboard de equipo. \\ \hline

GET & \path{/api/jugadores} & Obtiene el listado completo de jugadores. \\ \hline
GET & \path{/api/jugadores/mas-partidos} & Obtiene jugadores destacados por número de partidos. \\ \hline
GET & \path{/api/jugadores/:id_jugador} & Obtiene información de un jugador concreto. \\ \hline
GET & \path{/api/jugadores/:id_jugador/ratings} & Obtiene ratings y atributos de un jugador. \\ \hline
GET & \path{/api/jugadores/info/:id_jugador} & Obtiene el resumen del jugador. \\ \hline
GET & \path{/api/jugadores/partidos/:id_jugador} & Obtiene los partidos de un jugador. \\ \hline
GET & \path{/api/jugadores/trayectoria/:id_jugador} & Obtiene la trayectoria de un jugador. \\ \hline

GET & \path{/api/temporadas/annos} & Obtiene las temporadas disponibles. \\ \hline
GET & \path{/api/temporadas} & Obtiene la clasificación de una temporada. \\ \hline
GET & \path{/api/temporadas/partidos} & Obtiene los partidos de una temporada. \\ \hline
GET & \path{/api/temporadas/equipos} & Obtiene los equipos de una temporada. \\ \hline
GET & \path{/api/temporadas/info} & Obtiene el resumen de una temporada. \\ \hline
GET & \path{/api/temporadas/rankings} & Obtiene rankings de jugadores por temporada. \\ \hline
GET & \path{/api/temporadas/graficos} & Obtiene datos para gráficos de temporada. \\ \hline

GET & \path{/api/partidos} & Obtiene partidos filtrados por fecha, temporada o jornada. \\ \hline
GET & \path{/api/partidos/jornadas} & Obtiene jornadas disponibles para una temporada. \\ \hline
GET & \path{/api/partidos/h2h/:id1/:id2} & Obtiene el historial entre dos equipos. \\ \hline
GET & \path{/api/partidos/:id_partido/info} & Obtiene información general de un partido. \\ \hline
GET & \path{/api/partidos/:id_partido/eventos} & Obtiene eventos de un partido. \\ \hline
GET & \path{/api/partidos/:id_partido/alineaciones} & Obtiene alineaciones de un partido. \\ \hline
GET & \path{/api/partidos/:id_partido/stats_equipos} & Obtiene estadísticas de equipos en un partido. \\ \hline
GET & \path{/api/partidos/:id_partido/stats_jugadores} & Obtiene estadísticas de jugadores en un partido. \\ \hline

POST & \path{/api/usuarios/registro} & Registra un nuevo usuario. \\ \hline
POST & \path{/api/usuarios/login} & Inicia sesión. \\ \hline
GET & \path{/api/usuarios/:id} & Obtiene los datos de un usuario. \\ \hline
PUT & \path{/api/usuarios/:id} & Actualiza los datos de un usuario. \\ \hline
PUT & \path{/api/usuarios/:id/password} & Cambia la contraseña de un usuario. \\ \hline

GET & \path{/api/buscador?q=texto} & Realiza una búsqueda global de equipos y jugadores. \\ \hline
GET & \path{/api/favoritos/:id_usuario} & Obtiene los favoritos del usuario. \\ \hline
POST & \path{/api/favoritos/toggle} & Añade o elimina un favorito. \\ \hline
POST & \path{/api/chatbot} & Envía una pregunta al chatbot. \\ \hline

\end{longtable}

% \subsection{Inicio de sesión y registro}
% \subsection{Buscador global}
% El backend incluye un buscador global que permite localizar jugadores y equipos desde un único endpoint. Esta funcionalidad se implementa mediante la ruta /api/buscador, que recibe el texto de búsqueda a través del parámetro q.

% \begin{lstlisting}[caption = Ruta de buscador]
% const express = require('express');
% const router = express.Router();
% const buscadorController = require('../controllers/buscadorController');

% router.get('/', buscadorController.buscarGlobal);

% module.exports = router;
% \end{lstlisting}

% \begin{lstlisting}[caption=Endpoint de buscador]
% const buscarGlobal = async (req, res) => {
%   const { q } = req.query;
%   const busqueda = (q || '').trim();

%   if (!busqueda || busqueda.length < 2) {
%     return res.json([]);
%   }

%   const query = `
%     (SELECT id_jugador AS id, nombre AS principal, nacionalidad AS secundario,
%             foto AS imagen, 'jugador' AS tipo
%      FROM dim_jugador
%      WHERE LOWER(nombre) LIKE '%' || LOWER($1) || '%'
%      LIMIT 5)
%     UNION ALL
%     (SELECT id_equipo AS id, nombre_equipo AS principal, pais AS secundario,
%             logo AS imagen, 'equipo' AS tipo
%      FROM dim_equipo
%      WHERE LOWER(nombre_equipo) LIKE '%' || LOWER($1) || '%'
%      LIMIT 5);
%   `;

%   const result = await pool.query(query, [busqueda]);
%   res.json(result.rows);
% };
% \end{lstlisting}

% El buscador consulta de forma conjunta las tablas dim\_jugador y dim\_equipo,
% devolviendo los resultados con una estructura común: identificador, nombre
% principal, información secundaria, imagen y tipo de resultado. De esta forma,
% el frontend puede mostrar jugadores y equipos en una misma lista y navegar
% después al detalle correspondiente. Ademas, la consulta completa prioriza las coincidencias exactas frente a las coincidencias parciales, para que los resultados mas relevantes aparezcan antes.

% \subsection{Favoritos}

% El backend permite que cada usuario pueda marcar equipos y jugadores como
% favoritos. Esta funcionalidad se gestiona mediante el módulo favoritos, compuesto por una ruta para obtener los favoritos de un usuario y otra para alternar entre añadir o eliminar un favorito.

% \begin{lstlisting}
% const express = require('express');
% const router = express.Router();
% const favoritosController = require('../controllers/favoritosController');

% router.post('/toggle', favoritosController.toggleFavorito);
% router.get('/:id_usuario', favoritosController.getFavoritosUsuario);

% module.exports = router;
% \end{lstlisting}

% La función toggleFavorito recibe el identificador del usuario, el identificador del elemento marcado y el tipo de favorito, que puede ser equipo o jugador. Si el favorito ya existe, se elimina; si no existe, se inserta en la tabla h\_usuario\_favoritos.

% En el frontend, estos datos se consumen desde un contexto global de favoritos. Esto permite que las distintas pantallas puedan saber si un equipo o jugador ya está marcado como favorito y actualizar la interfaz al pulsar el botón con forma de corazón.

\subsubsection{Chatbot}
El \textit{backend} incorpora un modulo de chatbot que permite al usuario realizar
preguntas en lenguaje natural sobre los datos de LaLiga. La idea principal es transformar una pregunta escrita por el usuario en una consulta SQL segura, ejecutarla sobre la base de datos PostgreSQL y devolver una respuesta en lenguaje natural.

Este \textit{chatbot} hace uso de una APIKEY de Gemini proporcionada por Google AI Studio \cite{aistudio} lo que nos da permiso a usar sus modelos en su aplicación. En este caso haremos uso del modelo Gemini 2.5 Flash, el cual da un rendimiento tanto de precisión como de velocidad de respuesta suficiente para esta herramienta.

La conexión con la API de Gemini la configuramos en \path{config/gemini.js}, usando la APIKEY que hemos recibido y esta dentro de las variables de entorno.

En la ruta, tenemos un \textit{endpoint} \path{/api/chatbot/} que llama a la función del controlador \texttt{contestarPregunta}, y se encarga de todo el funcionamiento del \textit{chatbot}.

La petición recibe la pregunta del usuario en formato JSON y la recibe el controlador que coordina todo el flujo.

Una vez recibida la respuesta, podemos distinguir 4 pasos claves del controlador:
\begin{enumerate}
    \item Generación de SQL: La función generarSQL se encarga de construir el \textit{prompt} enviado a Gemini. Para ello combina la pregunta del usuario con un \textit{prompt} que describe el esquema disponible de la base de datos, las reglas de uso, las tablas permitidas y las condiciones que debe cumplir la consulta generada.
    \item Validación de seguridad: Esto es importante porque evita que el modelo genere operaciones que modifiquen o eliminen datos. Solo se aceptan consultas de lectura.
    \item Ejecución de consulta: se ejecuta la consulta en la base de datos y se devuelve un json con la respuesta obtenida, en caso contrario, se devuelve un mensaje de error.
    \item Generación de respuesta natural: Una vez obtenidos los datos, se hace una segunda llamada a Gemini para convertir el resultado SQL en una respuesta comprensible para el usuario.
\end{enumerate}

Esta implementación permite integrar inteligencia artificial en el \textit{backend} sin darle acceso directo e ilimitado a la base de datos. El modelo genera una consulta SQL, pero antes de ejecutarse pasa por una validación que restringe las operaciones permitidas. Posteriormente, los resultados se transforman en una respuesta breve y comprensible para el usuario, ocultando la complejidad de la consulta SQL.




%%%%%%%%%%%%%%%%%%%%%%%%%%%%%%%%%%%%%%%%%%%%%%%%%%%%%%%%%%%%%%%

